\begin{definition}[First-Main-Lemma error term]\label{mod:delta-errorterm}
Let $K/F$ be a finite extension of nonarchimedean local fields for which
$S(K/F)=\texttt{NormCharacter}(F,K)$ is finite, and let $\Delta_F,\Delta_K$
be proposed local-constant functions.  Lean defines the error term with the
manuscript's exact left-over-right orientation:
\[
 \texttt{errorTerm}(\chi_F,\psi_F)
 =\frac{\texttt{firstMainLeftSide}(\chi_F,\psi_F)}
        {\texttt{firstMainRightSide}(\chi_F,\psi_F)}
 =\frac{
  \Delta_K(\chi_F\circ N_{K/F},\psi_F\circ\operatorname{Tr}_{K/F})
  \prod_{\mu\in S(K/F)}\Delta_F(\mu,\psi_F)}
 {\prod_{\mu\in S(K/F)}\Delta_F(\mu\chi_F,\psi_F)}.
\]

The predicate \texttt{IsDeltaFiniteLocalConstant} says that a proposed
local-constant function agrees with \texttt{deltaFinite} whenever exact
multiplicative and additive conductor packages and an admissible denominator
are supplied.  It is a computational compatibility condition, not a
nonvanishing assumption.  From it, Lean derives nonvanishing directly from
\texttt{deltaFinite\_ne\_zero}.

For fixed $\chi_F,\psi_F$, the structure
\texttt{FirstMainComputationalData} packages exact conductor data for the
extension-field pullbacks and for every base-field character $\mu$ and
$\mu\chi_F$ occurring in the two products.  Both families are indexed by the
full type $S(K/F)$, including its trivial character.  With this data Lean
first proves every denominator factor nonzero and hence proves
\[
 \texttt{firstMainRightSide}(\chi_F,\psi_F)\ne0.
\]
It separately proves the extension-field factor and every factor of the
remaining numerator product nonzero before proving the numerator and then
the error term nonzero.  No product is cancelled before its nonvanishing is
established.

The proved rearrangement API includes both
\[
 \texttt{errorTerm}\cdot\texttt{firstMainRightSide}
   =\texttt{firstMainLeftSide},\qquad
 \texttt{firstMainRightSide}\cdot\texttt{errorTerm}
   =\texttt{firstMainLeftSide},
\]
and, for every $z\in\mathbf C$, the two equivalent numerator/denominator
forms of \texttt{errorTerm}$=z$.  At $z=1$ this gives both implications of
\[
 \texttt{errorTerm}(\chi_F,\psi_F)=1
 \quad\Longleftrightarrow\quad
 \texttt{FirstMainIdentity}(\chi_F,\psi_F).
\]
The corresponding universally quantified equivalence is exported for the
project's exact proposition \texttt{FirstMainStatement}.  These equivalences
do not assume the First Main identity.

Lean also exports the generic computational bridge
\texttt{firstMainIdentity\_of\_deltaFinite\_product}.  Given computational
specifications for $\Delta_F$ and $\Delta_K$, exact conductor packages for
the extension pullbacks and every $\mu$ and $\mu\chi_F$, admissible
denominators $\gamma_K$, $\gamma_{\mu}$, and $\gamma_{\mu\chi}$, it takes the
single explicit equality
\[
 \Delta^{\mathrm{fin}}_K(\chi_K,\psi_K;\gamma_K)
 \prod_{\mu\in S(K/F)}
   \Delta^{\mathrm{fin}}_F(\mu,\psi_F;\gamma_{\mu})
 =
 \prod_{\mu\in S(K/F)}
   \Delta^{\mathrm{fin}}_F(\mu\chi_F,\psi_F;\gamma_{\mu\chi})
\]
and deduces \texttt{FirstMainIdentity}.  Each displayed product is over the
full norm-character type, including the identity character.  The proof
rewrites the proposed local constants through their \texttt{deltaFinite}
specifications, uses the displayed equality with its left/right orientation
unchanged, proves the denominator product nonzero through the existing
nonvanishing API, and then applies
\texttt{errorTerm}$=1\Longleftrightarrow$\texttt{FirstMainIdentity}.

For twist invariance, Lean constructs multiplication by a fixed
$\nu\in S(K/F)$ directly on the exact \texttt{NormCharacter} subtype,
together with the inverse multiplication by $\nu^{-1}$.  This gives a
genuine equivalence of the full norm-character type and the reindexing
identity
\[
 \prod_{\mu\in S(K/F)} f(\nu\mu)
 =\prod_{\mu\in S(K/F)} f(\mu),
\]
again retaining the trivial-character factor.  The defining property
$\nu\circ N_{K/F}=1$ proves that the extension-field pullback of
$\nu\chi_F$ equals that of $\chi_F$, while the displayed bijection reindexes
the entire denominator.  Therefore Lean proves
\[
 \texttt{errorTerm}(\nu\chi_F,\psi_F)
 =\texttt{errorTerm}(\chi_F,\psi_F)
\]
without assuming a First Main identity or any extra nonvanishing hypothesis.
\LeanFile{LanglandsFirstMainLemma/Delta/ErrorTerm.lean}
\lean{LanglandsFirstMainLemma.IsDeltaFiniteLocalConstant}
\lean{LanglandsFirstMainLemma.IsDeltaFiniteLocalConstant.apply_ne_zero}
\lean{LanglandsFirstMainLemma.FirstMainComputationalData}
\lean{LanglandsFirstMainLemma.errorTerm_denominator_ne_zero}
\lean{LanglandsFirstMainLemma.errorTerm_numerator_ne_zero}
\lean{LanglandsFirstMainLemma.errorTerm}
\lean{LanglandsFirstMainLemma.errorTerm_ne_zero}
\lean{LanglandsFirstMainLemma.errorTerm_mul_denominator}
\lean{LanglandsFirstMainLemma.denominator_mul_errorTerm}
\lean{LanglandsFirstMainLemma.errorTerm_eq_iff_numerator_eq_mul_denominator}
\lean{LanglandsFirstMainLemma.errorTerm_eq_iff_denominator_mul_eq_numerator}
\lean{LanglandsFirstMainLemma.errorTerm_eq_one_iff_firstMainIdentity}
\lean{LanglandsFirstMainLemma.firstMainIdentity_of_deltaFinite_product}
\lean{LanglandsFirstMainLemma.normCharacterLeftMul}
\lean{LanglandsFirstMainLemma.normCharacterLeftMulEquiv}
\lean{LanglandsFirstMainLemma.normCharacter_fullProduct_leftMul}
\lean{LanglandsFirstMainLemma.normQuasiChar_normCharacter_mul}
\lean{LanglandsFirstMainLemma.firstMainRightSide_normCharacter_mul}
\lean{LanglandsFirstMainLemma.firstMainLeftSide_normCharacter_mul}
\lean{LanglandsFirstMainLemma.errorTerm_normCharacter_mul}
\lean{LanglandsFirstMainLemma.firstMainStatement_iff_errorTerm_eq_one}
\leanok
\SourceText{Definition eq:error-term and Lemma lem:twist-invariance.}
\uses{mod:delta-firstmainstatement,mod:delta-nonvanishing}
\FormalizationContract The definition is exactly the quotient of the two
previously defined First Main sides, with numerator on the left and
denominator on the right.  Nonvanishing hypotheses are never accepted as
inputs: every required nonzero local-constant value is derived from exact
conductor data, computational compatibility, and
\texttt{deltaFinite\_ne\_zero}.  The full-product permutation is an explicit
equivalence of \texttt{NormCharacter}, not a reindexing assertion without a
bijection.  All products include the trivial norm character.  The fixed-point
and statement equivalences prove both directions and assume no First Main
identity.  The generic computational-product bridge accepts exactly the
extension factor and two full norm-character products, preserves their
left/right orientation, and performs no cancellation before the existing
denominator nonvanishing theorem has been applied.
\end{definition}
